\documentclass[11pt, letterpaper]{article}

\input{/home/hyperion/.config/LatexHub/preambles/base.tex}
\input{/home/hyperion/.config/LatexHub/preambles/tikz.tex}
\input{/home/hyperion/.config/LatexHub/preambles/diagrams.tex}
\input{/home/hyperion/.config/LatexHub/preambles/macros-cat-theory.tex}
\input{/home/hyperion/.config/LatexHub/preambles/macros-algebra.tex}
\input{/home/hyperion/.config/LatexHub/preambles/homework-formatting-global.tex}
\input{/home/hyperion/.config/LatexHub/preambles/refs-basic.tex}
\crefname{problem}{problem}{problems}
\Crefname{problem}{Problem}{Problems}
\usepackage{titlepageBU}
\title{Assignment 01}
\courseID{MA511}
\begin{document}
\makeproblem

\begin{notation}
	We write $\mathbb{N} $ for the set of naturals, which we only use for convenience of notation, since sets have not been formally introduced. Given $n\in\mathbb{N} $, we write $s(n)$ for the successor of $n$. Multiplication is always denoted by juxtaposition or $\cdot $. References to in-class results/definitions are either given by name, or cited in the first lecture notes \cite{lec-1}.
\end{notation}


\begin{problem}[Addition is commutative]\label{prob:1}
	For any naturals $n,m$, we have $n+m=m+n$.
\end{problem}
\begin{proof}
	We proceed by induction on $n$. Fix a natural $m$, and let $n=0$. By definition of addition, we have $0+m=m$, and by \cite{lec-1} \S2.2 Lemma 1, we have $m+0=m$. Thus, $m+0=m=0+m$.

	Now assume $n+m=m+n$ for some natural $n$, and consider the expression $s(n)+m$. By definition of addition, $s(n)+m = s(n+m)$. By the inductive hypothesis, $n+m=m+n$, and thus $s(n+m) = s(m+n)$. By \cite{lec-1} \S2.2 Lemma 2, we have $s(m+n) = m+s(n)$. Thus, we have equalities
	\[
		s(n)+m = s(n+m) = s(m+n) = m+s(n).
	\]
	The statement now follows by induction.
\end{proof}

\begin{problem}[Cancellation law]\label{prob:2}
	Let $a,b,c$ be naturals such that $a+b=a+c$. Then $b=c$.
\end{problem}
\begin{proof}
	We proceed by induction on $a$. Let $a=0$, and let $b,c\in\mathbb{N} $ such that $a+b=a+c$. Then by definition of addition, $0+b=b$ and $0+c=c$. Thus,
	\[
		b=0+b=0+c=c
	.\]
	Now assume $a+b=a+c$ implies that $b=c$ for some natural $a$ and all naturals $b,c$, and suppose that $s(a)+b=s(a)+c$ for some $b,c\in\mathbb{N} $. Then by definition of addition, we have
	\[
		s(a+b)=s(a)+b = s(a)+c = s(a+c)
	.\]
	By Peano axiom 4, it follows that $a+b=a+c$. By the inductive hypothesis, $b=c$. By the method of induction, we are done.
\end{proof}

\begin{problem}[Order is transitive]\label{prob:3}
	Prove that ordering is transitive.
\end{problem}
\begin{proof}
	Let $a,b,c$ be naturals with $c\leq b\leq a$. We must show $c\leq a$. The following lemma is required.

	\begin{lemma}\label{lma:3-1}
		Let $a,b,c$ be naturals. Then $(a+b)+c=a+(b+c)$.
	\end{lemma}
	\begin{proof}
		We proceed by induction on $a$. Fix naturals $b,c$, and let $a=0$, and recall that $0+b=b$ and $0+(b+c)=b+c$ by definition of addition. Thus
		\[
			(0+b)+c = b+c = 0+(b+c)
		.\]
		Now let $a$ be an arbitrary natural, and consider the successor $s(a)$. We have
		\begin{align*}
			(s(a)+b)+c &=s(a+b)+c &\text{Definition of addition}\\
				   &=s((a+b)+c)&\text{Definition of addition}\\
				   &=s(a+(b+c))&\text{Inductive hypothesis}\\
				   &=s(a)+(b+c)&\text{Definition of addition}
		.\end{align*}
		By the method of induction, the statement is proven.
	\end{proof}
	We now continue with \cref{prob:3}. Since $c\leq b$, there is a natural $n$ such that $c+n=b$. Since $b\leq a$, there is a natural $m$ such that $b+m=a$. We have
	\[
		a=b+m = (c+n)+m \overset{!}{=}c+(n+m)
	,\]
	where the equality marked with a $!$ follows by \cref{lma:3-1}. Since $n+m$ is a natural, by definition of $\leq $, we have $c\leq a$.
\end{proof}

\begin{problem}[Multiplication is commutative]\label{prob:4}
	Let $n,m$ be naturals. Then $nm=mn$.
\end{problem}
\begin{proof}
	We use $ab=ba$ instead of $nm=mn$, to improve readability. We proceed by induction on $a$. Let $a=0$, and by definition of multiplication, $0b=0$ for all $b\in\mathbb{N} $. We must show $b0=0$. We proceed by induction on $b$. When $b=0$, we observe that $0\cdot 0=0$ is true by definition of multiplication. Now assume $b0=0$, and consider the expression $s(b)0$. By definition of multiplication, $s(b)0=(b0)+0$. By the inductive hypothesis, $b0=0$, and by definition of addition, $0+0=0$. Thus,
	\[
		s(b)0=(b0)+0 = 0+0=0
	.\]
	Thus, $b0=0$ holds for all $b$.

	Now we assume the inductive hypothesis for $a$; that is, we assume $ab=ba$ for some arbitrary natural $a$ and all naturals $b$. By definition of multiplication, 
	\[
		s(a)b=ab+b = ba+b 
	.\]
	It suffices to show $bs(a)=ba+b$ for all naturals $b$. We proceed by induction on $b$. First, take $b=0$. Notice that $0s(a) = 0$ by definition of multiplication, and $0a + 0 = 0+0=0$ by definition of multiplication and addition. Now assume the statement holds for $b$ an arbitrary natural. We want to show $s(b)s(a) = s(b)a+s(b)$. We have
	\begin{align*}
		s(b)s(a) &=bs(a)+s(a)&\text{Definition of multiplication}\\
			 &=(ba+b)+s(a)&\text{Inductive hypothesis} \\
			 &=ba+(b+s(a))&\text{\Cref{lma:3-1}}\\
			 &=ba+(s(a)+b)&\text{\Cref{prob:1}}\\
			 &=ba+s(a+b)&\text{Definition of addition} \\
			 &=ba+s(b+a)&\text{\Cref{prob:1}} \\
			 &=ba+(s(b)+a)&\text{Definition of addition}  \\
			 &=ba+(a+s(b))&\text{\Cref{prob:1}} \\
			 &=(ba+a)+s(b)&\text{\Cref{lma:3-1}}\\
			 &=s(b)a+s(b)&\text{Definition of multiplication} 
	.\end{align*}
	Thus, the statement is proven.
\end{proof}

\begin{problem}[$\mathbb{N} $ has no zero-divisors]\label{prob:5}
	Let $n,m$ be naturals. Then $nm=0$ if and only if at least one of $n,m$ is $0$.
\end{problem}
We once again use $a,b$ in place of $n,m$ to improve readability. The proof of one direction requires a lemma, which was proven in class but is not contained in the notes.
\begin{lemma}\label{lma:5-1}
	Let $a$ be a natural. Then $a$ is either 0, or a successor.
\end{lemma}
\begin{proof}
	We proceed by induction. If $a=0$, then $a=0$, and by Peano axiom 2 $a$ is not a successor.
		Now let $a$ be an arbitrary natural. Then $s(a)$ is a successor. Moreover, by Peano axiom 2, $s(a)\neq 0$. The statement is now proven.
\end{proof}
\begin{proof}[Proof of \cref{prob:5}]
	($\implies $) Let $ab=0$. If $a=0$ or $b=0$ we are done, so suppose for contradiction that $a,b\neq 0$. Applying \cref{lma:5-1} to $a,b\neq 0$, we see that $a=s(\widetilde{a} )$ and $b=s(\widetilde{b} )$ for some naturals $\widetilde{a} ,\widetilde{b}$. By the definition of multiplication, we compute
	\[
		ab = s(\widetilde{a} )s(\widetilde{b} ) = \widetilde{a} s(\widetilde{b} )+s(\widetilde{b})
	.\]
	Applying \cref{prob:4} and the definition of multiplication, we have
	\[
		\widetilde{a} s(\widetilde{b} )+s(\widetilde{b})  = s(\widetilde{b}) \widetilde{a} +s(\widetilde{b}) =(\widetilde{b} \widetilde{a} +\widetilde{a} )+s(\widetilde{b} )
	.\]
	We now have
	\begin{align*}
		(\widetilde{b} \widetilde{a} +\widetilde{a} )+s(\widetilde{b} )&=s(\widetilde{b} )+( \widetilde{b} \widetilde{a} +\widetilde{a} ) &\text{\Cref{prob:1}} \\
											     &=s( \widetilde{b} +(\widetilde{b} \widetilde{a} +\widetilde{a} ))  &\text{Definition of addition}
	.\end{align*}
	Notice that the equation is thus equal to a successor, and by Peano axiom 2, all successors are nonzero. Thus,
	\[
		ab=s( \widetilde{b} +(\widetilde{b} \widetilde{a} +\widetilde{a} )) \neq 0
	,\]
	and the first direction is proven.
	

	($\impliedby $) Let at least one of $a,b$ be $0$. By \cref{prob:4}, we can take $a=0$ without loss of generality. Then by definition of multiplication,
	\[
		ab=0b=0
	.\]
	The statement is proven.
\end{proof}

\begin{problem}
	For any naturals $a,b,c$, we have
	\begin{equation}\label{eq:6-1}
		a(b+c)=ab+ac
	\end{equation}
	and
	\begin{equation}\label{eq:6-2}
		(b+c)a=ba+ca.
	\end{equation}
\end{problem}
\begin{proof}
	Assume \eqref{eq:6-1} holds. Then \eqref{eq:6-2} follows by \cref{prob:4}:
	\[
		(b+c)a=a(b+c)=ab+ac=ba+ca
	.\]
	Thus, it suffices to prove \eqref{eq:6-1}. We proceed by induction on $a$. If $a=0$, then $0(b+c)=0$ by definition of multiplication. Additionally, $0b=0c=0$ by definition of multiplication, and thus $0b+0c=0+0=0$ by definition of addition. Thus, $0(b+c)=0b+0c$.

	Now assume $a(b+c)=ab+ac$ for some arbitrary natural $a$. We have
	\begin{align*}
		s(a)(b+c)&=a(b+c)+(b+c)&\text{Definition of multiplication}\\
			 &=(ab+ac)+(b+c)&\text{Inductive hypothesis} \\
			 &=(ab+b)+(ac+c)&\text{Repeated applications of \cref{prob:1,lma:3-1}}\\
			 &=s(a)b+s(a)c&\text{Definition of multiplication}
	.\end{align*}
	By the method of induction, we are done.
\end{proof}





\bibliographystyle{alpha}
\bibliography{refs-1.bib}
\end{document}
